Vision-language models (VLMs) are increasingly used as judges that pick the best of several generated images, so their choices decide what users see. Such judges are usually validated by how well their scores agree with human ratings, not by the images they choose. We audit VLM judges as decision-makers: on 300 culturally situated prompts, we compare the image a judge returns with human ratings it never sees and with random choice from the same candidates, and we repeat every decision with the candidates reordered. A 4B-parameter judge barely beats random selection and falls short of a CLIP similarity baseline. Its choices show strong position bias: it picks the first image shown in 49\% of calls, where chance would give 28\%, and reordering changes the image it returns on 60\% of prompts. For this judge, agreement across orders is informative: decisions that survive reordering are much better than random, whereas agreement with a weaker second judge keeps the wrong ones. An 8B judge from the same family shows almost no position bias and beats the CLIP baseline, and for it the same order-agreement filter mostly discards good decisions. Agreement helps only when it targets the judge's failure mode, so such filters must be re-audited whenever the judge changes. The 4B judge's slight rise in stereotype ratings is no longer detectable after aggregating across orders or with the larger judge.
